\section{Alias et chevauchement mémoire}
\label{secAliasChevauchement}

Avant de pouvoir convertir quoi que ce soit, il faut d'abord pouvoir
\textbf{nommer} la zone mémoire brute échangée avec le réseau. C'est le
rôle de l'\textbf{alias}.

\begin{UPSTIinfor}{Alias et chevauchement mémoire}
Sous Control Expert, une variable peut être déclarée par
\textbf{chevauchement} (\emph{overlap}) d'une autre : elle occupe alors
exactement le même espace mémoire que la variable qu'elle chevauche, sans
qu'aucune copie ne soit effectuée. On parle d'\textbf{alias} parce que les
deux variables désignent le même contenu physique sous deux noms et,
éventuellement, deux types différents.

Cette technique est indispensable chaque fois qu'un coupleur réseau
(comme le \texttt{BMX~NOC~0401.2}) ne connaît que des zones \texttt{\%MW}
génériques, numérotées, alors que le programme applicatif a besoin de
noms et de types porteurs de sens.
\end{UPSTIinfor}

\textbf{Exemple.} Le coupleur \texttt{NOC\_ETHIP\_CONVOYEUR} reçoit du
PFC200 un octet de données à l'adresse \texttt{\%MW10} (poids faible). Si
le quatrième bit de cet octet correspond au capteur « palette présente »,
on déclare une variable \texttt{m\_nixPalettePresente} de type
\texttt{BOOL}, chevauchant précisément le bit~3 de \texttt{\%MW10} :

\begin{lstlisting}
m_nixPalettePresente AT %MW10.3 : BOOL;
\end{lstlisting}

Aucune conversion n'est nécessaire ici : le bit \emph{est} directement
l'information recherchée, l'alias suffit à lui donner un nom exploitable
(c'est le cas « à un seul niveau » du PFC200, déjà rencontré en TP1).
C'est précisément lorsque cela ne suffit plus -- lorsque l'information
recherchée est encodée sur plusieurs bits, répartie, ou calculée à partir
d'octets bruts -- qu'un bloc fonctionnel de conversion devient nécessaire,
comme on le verra dans les sections suivantes pour le M221 et pour la baie
CS9.

\begin{UPSTIactivite}[4][Chevauchement sur Schneider M340][][][TD]
\UPSTIquestion{Le coupleur \texttt{NOC\_ETHIP\_CONVOYEUR} reçoit du
PFC200, à l'adresse \texttt{\%MW11}, un octet dont le bit~0 indique l'état
du bouton d'arrêt d'urgence du poste guichet (\texttt{TRUE} = défaut). En
complétant la déclaration ci-dessous, créer un alias
\texttt{m\_nixArretUrgenceGuichet} chevauchant ce bit.}

\begin{lstlisting}
m_nixArretUrgenceGuichet AT (*@\trou{2.2cm}@*) : (*@\trou{1.4cm}@*);
\end{lstlisting}

\UPSTIquestion{Le même octet comporte, sur ses bits~1 et~2, un code sur
2~bits indiquant la couleur de la palette détectée (\texttt{00} = absente,
\texttt{01} = rouge, \texttt{10} = verte, \texttt{11} = bleue). Expliquer
en une phrase pourquoi un simple alias sur ces deux bits ne permettrait
\textbf{pas}, à lui seul, d'obtenir directement une variable
\texttt{m\_couleurPalette} de type énuméré portant les noms
\texttt{ROUGE}/\texttt{VERTE}/\texttt{BLEUE}/\texttt{ABSENTE} : quelle
étape supplémentaire serait alors nécessaire ?}
\end{UPSTIactivite}
